%! AP Statistics — Unit 2, Lesson 2.2 — Slides (Beamer)
\documentclass[aspectratio=169, 11pt]{beamer}
\usepackage{apstats-beamer}

%=============================================================================
\begin{document}

% ─────────────────────────────────────────────────────────────────────────────
% SLIDE 1 — LESSON TITLE
% ─────────────────────────────────────────────────────────────────────────────
{
\setbeamercolor{background canvas}{bg=navy}
\begin{frame}[plain]
  \begin{tikzpicture}[remember picture, overlay]
    \fill[navylight] (current page.north west)
      rectangle ([yshift=-1.35cm]current page.north east);
  \end{tikzpicture}
  \vspace{2.8cm}
  \hspace{0.55cm}%
  \begin{minipage}{0.9\textwidth}
    {\color{goldacc}\bfseries\small LESSON 2.2}\\[0.5cm]
    {\color{white}\bfseries\LARGE Representing Two\\[0.25cm]
      Categorical Variables}\\[1.4cm]
    {\color{skymid}\small AP Statistics~$\cdot$~Unit 2: Exploring Two-Variable Data}
  \end{minipage}
\end{frame}
}

% ─────────────────────────────────────────────────────────────────────────────
% SLIDE 2 — PRIMARY OBJECTIVE
% ─────────────────────────────────────────────────────────────────────────────
{
\setbeamercolor{background canvas}{bg=sky}
\begin{frame}[t, plain]
  \begin{tikzpicture}[remember picture, overlay]
    \fill[navy] (current page.north west)
      rectangle ([yshift=-0.25cm]current page.north east);
  \end{tikzpicture}
  \vspace{0.45cm}\par
  \hspace{0.55cm}%
  \begin{minipage}{0.9\textwidth}

    \sectionlabel{Primary Objective}

    \normalsize
    Students will construct and interpret two-way tables and segmented bar charts
    for two categorical variables, and use \textbf{conditional relative frequencies}
    to determine whether the variables are associated (Big Idea: UNC, Skill 2.D).

    \vspace{0.6cm}

    \normalsize\textbf{\color{navy}By the end of this lesson, I will be able to:}
    \smallskip
    \begin{itemize}
      \setlength{\itemsep}{0.3em}
      \item Construct a two-way table and identify joint and marginal frequencies.
      \item Compute joint and conditional relative frequencies and explain the difference.
      \item Construct and interpret a segmented bar chart to determine association.
    \end{itemize}

  \end{minipage}
\end{frame}
}

% ─────────────────────────────────────────────────────────────────────────────
% SLIDE 3 — HOOK
% ─────────────────────────────────────────────────────────────────────────────
{
\setbeamercolor{background canvas}{bg=hookbg}
\setbeamercolor{itemize item}{fg=white}
\begin{frame}[t, plain]
  \vspace{0.45cm}\par
  \hspace{0.55cm}%
  \begin{minipage}{0.9\textwidth}

    {\color{white}\bfseries\footnotesize\MakeUppercase{Hook}}

    \vspace{0.3cm}
    {\color{white}\large\bfseries A researcher surveys 240 students and records:}\\[0.3cm]
    {\color{skymid} Variable 1: Grade level --- Underclassman or Upperclassman}\\[0.15cm]
    {\color{skymid} Variable 2: Preferred lunch --- Cafeteria, Off-campus, or Home}

    \vspace{0.5cm}
    {\color{goldacc}\bfseries Statistical question:}\\[0.2cm]
    {\color{white}\itshape ``Is there an association between grade level and lunch preference?''}

    \vspace{0.5cm}
    {\color{white}\bfseries Discussion:}
    \begin{itemize}
      \setlength{\itemsep}{0.35em}
      \item {\color{white}Why can't a single bar chart answer this question?}
      \item {\color{white}What kind of table or chart would let you compare lunch
        preferences \emph{across} grade levels?}
    \end{itemize}

  \end{minipage}
\end{frame}
}

% ─────────────────────────────────────────────────────────────────────────────
% SLIDE 4 — BUILDING THE TWO-WAY TABLE
% ─────────────────────────────────────────────────────────────────────────────
{
\setbeamercolor{background canvas}{bg=white}
\begin{frame}[t, plain]
  \navyheader{Step 1 --- Build the Two-Way Table (UNC-1.P.3)}
  \hspace{0.4cm}%
  \begin{minipage}{0.93\textwidth}
    \vspace{0.3cm}
    \small
    \renewcommand{\arraystretch}{1.8}
    \begin{tabular}{l cccc}
    \rowcolor{sky}
    \textbf{Grade Level} & \textbf{Cafeteria} & \textbf{Off-campus} & \textbf{Home} & \textbf{Row Total} \\
    \hline
    Underclassman & 54 & 18 & 48 & \textcolor{navy!60}{\textit{fill in}} \\
    Upperclassman & 36 & 42 & 42 & \textcolor{navy!60}{\textit{fill in}} \\
    \hline
    \textbf{Col.\ Total} & \textcolor{navy!60}{\textit{fill}} & \textcolor{navy!60}{\textit{fill}} & \textcolor{navy!60}{\textit{fill}} & \textcolor{navy!60}{\textit{fill}} \\
    \end{tabular}

    \vspace{0.5cm}
    \begin{columns}[T]
      \begin{column}{0.48\textwidth}
        \textbf{\color{navy}Vocabulary:}
        \begin{itemize}\setlength{\itemsep}{0.3em}
          \item \textbf{Joint frequency} --- count in an interior cell (describes both variables)
          \item \textbf{Marginal frequency} --- row or column total (describes one variable)
        \end{itemize}
      \end{column}
      \begin{column}{0.48\textwidth}
        \textbf{\color{navy}Grand total:}\\
        Sum of row totals = sum of column totals = \textbf{240}
      \end{column}
    \end{columns}
  \end{minipage}
\end{frame}
}

% ─────────────────────────────────────────────────────────────────────────────
% SLIDE 5 — JOINT vs. CONDITIONAL RELATIVE FREQUENCIES
% ─────────────────────────────────────────────────────────────────────────────
{
\setbeamercolor{background canvas}{bg=greenbg}
\begin{frame}[t, plain]
  \begin{tikzpicture}[remember picture, overlay]
    \fill[greenacc] (current page.north west)
      rectangle ([yshift=-1.35cm]current page.north east);
    \node[anchor=west, text=white, font=\bfseries\large, inner xsep=0.55cm]
      at ([yshift=-0.68cm]current page.north west)
      {Step 2 --- Joint vs.\ Conditional Relative Frequencies (UNC-1.P.4)};
  \end{tikzpicture}
  \vspace{1.05cm}\par
  \hspace{0.4cm}%
  \begin{minipage}{0.93\textwidth}
    \vspace{1em}
    \small
    \begin{columns}[T]
      \begin{column}{0.48\textwidth}
        \textbf{\color{greenacc}Joint relative frequency}\\[4pt]
        $\dfrac{\text{cell count}}{\text{grand total}}$\\[8pt]
        Example: underclassmen who eat in the cafeteria\\[4pt]
        $\dfrac{54}{240} = 0.225 = 22.5\%$ of \textbf{all students}
      \end{column}
      \begin{column}{0.48\textwidth}
        \textbf{\color{greenacc}Conditional relative frequency}\\[4pt]
        $\dfrac{\text{cell count}}{\text{row (or column) total}}$\\[8pt]
        Example: among underclassmen, proportion who eat in the cafeteria\\[4pt]
        $\dfrac{54}{120} = 0.45 = 45\%$ of \textbf{underclassmen}
      \end{column}
    \end{columns}
    \vspace{0.6cm}
    \textbf{\color{navy}Key question:} To detect association, which do we use and why?\\
    \textit{Use \textbf{conditional} --- it holds group size constant (``out of each group'').}
  \end{minipage}
\end{frame}
}

% ─────────────────────────────────────────────────────────────────────────────
% SLIDE 6 — CONDITIONAL DISTRIBUTION TABLE
% ─────────────────────────────────────────────────────────────────────────────
{
\setbeamercolor{background canvas}{bg=white}
\begin{frame}[t, plain]
  \navyheader{Step 3 --- Complete the Conditional Distribution Table}
  \hspace{0.4cm}%
  \begin{minipage}{0.93\textwidth}
    \vspace{0.3cm}
    \small
    Divide each row by its row total; round to the nearest percent.

    \vspace{0.3cm}
    \renewcommand{\arraystretch}{1.8}
    \begin{tabular}{l cccc}
    \rowcolor{sky}
    \textbf{Grade Level} & \textbf{Cafeteria} & \textbf{Off-campus} & \textbf{Home} & \textbf{Total} \\
    \hline
    Underclassman & 45\% & 15\% & 40\% & 100\% \\
    Upperclassman & 30\% & 35\% & 35\% & 100\% \\
    \end{tabular}

    \vspace{0.5cm}
    \textbf{Compare the rows:}
    \begin{itemize}\setlength{\itemsep}{0.3em}
      \item Underclassmen: mostly cafeteria (45\%), very few go off-campus (15\%)
      \item Upperclassmen: cafeteria use drops to 30\%, off-campus jumps to 35\%
      \item The rows are \textbf{different} $\to$ evidence of an \textbf{association}
    \end{itemize}
  \end{minipage}
\end{frame}
}

% ─────────────────────────────────────────────────────────────────────────────
% SLIDE 7 — SEGMENTED BAR CHART
% ─────────────────────────────────────────────────────────────────────────────
{
\setbeamercolor{background canvas}{bg=white}
\begin{frame}[t, plain]
  \navyheader{Step 4 --- Segmented Bar Chart (UNC-1.P.1, UNC-1.P.2)}
  \hspace{0.4cm}%
  \begin{minipage}{0.93\textwidth}
    \vspace{0.3cm}
    \begin{columns}[T]
      \begin{column}{0.48\textwidth}
        \small
        \textbf{\color{navy}Construction rules:}
        \begin{enumerate}\setlength{\itemsep}{0.3em}
          \item One bar per group (explanatory variable).
          \item Each bar reaches 100\%.
          \item Divide bar into segments using \textbf{conditional} relative frequencies.
          \item Label segments; use a consistent color legend.
        \end{enumerate}
        \vspace{6pt}
        \textbf{\color{navy}Reading the chart:}\\
        \textit{If bars look the same $\to$ no association.}\\
        \textit{If bars look different $\to$ association.}
      \end{column}
      \begin{column}{0.48\textwidth}
        \begin{tikzpicture}[x=1.5cm, y=0.033cm]
          \draw[thick,->] (0,0) -- (3.0,0);
          \draw[thick,->] (0,0) -- (0,115) node[above,font=\tiny] {\%};
          \foreach \y in {0,25,50,75,100}{
            \draw (0,\y) -- (-0.07,\y) node[left,font=\tiny] {\y};
          }
          \node[below,font=\small] at (0.75,0) {Under};
          \node[below,font=\small] at (2.25,0) {Upper};
          \fill[navy!70]     (0.2,  0) rectangle (1.3, 45);
          \fill[sky!80!navy] (0.2, 45) rectangle (1.3, 60);
          \fill[greenbg!90]  (0.2, 60) rectangle (1.3,100);
          \draw (0.2,0) rectangle (1.3,100);
          \fill[navy!70]     (1.7,  0) rectangle (2.8, 30);
          \fill[sky!80!navy] (1.7, 30) rectangle (2.8, 65);
          \fill[greenbg!90]  (1.7, 65) rectangle (2.8,100);
          \draw (1.7,0) rectangle (2.8,100);
          \fill[navy!70]   (-0.05,108) rectangle (0.13,116); \node[right,font=\tiny] at (0.15,112) {Cafe};
          \fill[sky!80!navy] (0.62,108) rectangle (0.80,116); \node[right,font=\tiny] at (0.82,112) {Off};
          \fill[greenbg!90] (1.15,108) rectangle (1.33,116); \node[right,font=\tiny] at (1.35,112) {Home};
        \end{tikzpicture}
      \end{column}
    \end{columns}
  \end{minipage}
\end{frame}
}

% ─────────────────────────────────────────────────────────────────────────────
% SLIDE 8 — VOCABULARY
% ─────────────────────────────────────────────────────────────────────────────
{
\setbeamercolor{background canvas}{bg=greenbg}
\begin{frame}[t, plain]
  \begin{tikzpicture}[remember picture, overlay]
    \fill[greenacc] (current page.north west)
      rectangle ([yshift=-1.35cm]current page.north east);
    \node[anchor=west, text=white, font=\bfseries\large, inner xsep=0.55cm]
      at ([yshift=-0.68cm]current page.north west)
      {Vocabulary \& Key Concepts};
  \end{tikzpicture}
  \vspace{1.05cm}\par
  \hspace{0.4cm}%
  \begin{minipage}{0.93\textwidth}
    \vspace{1em}
    \footnotesize
    \begin{multicols}{2}
      \textbf{\color{greenacc}Two-way table} --- organizes two categorical variables simultaneously; rows = one variable, columns = other\\[4pt]
      \textbf{\color{greenacc}Joint frequency} --- count in a single interior cell; describes both variables at once\\[4pt]
      \textbf{\color{greenacc}Marginal frequency} --- row or column total; one variable, ignoring the other\\[4pt]

      \columnbreak

      \textbf{\color{greenacc}Joint relative frequency} --- cell count $\div$ grand total (UNC-1.P.4)\\[4pt]
      \textbf{\color{greenacc}Conditional relative frequency} --- cell count $\div$ row or column total; holds one variable fixed\\[4pt]
      \textbf{\color{greenacc}Association} --- present when conditional distributions \emph{differ} across groups
    \end{multicols}
  \end{minipage}
\end{frame}
}

% ─────────────────────────────────────────────────────────────────────────────
% SLIDE 9 — GROUP ACTIVITY
% ─────────────────────────────────────────────────────────────────────────────
{
\setbeamercolor{background canvas}{bg=redbg}
\setbeamercolor{enumerate item}{fg=redacc}
\begin{frame}[t, plain]
  \begin{tikzpicture}[remember picture, overlay]
    \fill[redacc] (current page.north west)
      rectangle ([yshift=-1.35cm]current page.north east);
    \node[anchor=west, text=white, font=\bfseries\large,
          inner xsep=0.55cm, inner ysep=0pt]
      at ([yshift=-0.675cm]current page.north west)
      {Group Activity};
  \end{tikzpicture}
  \vspace{1.05cm}\par
  \hspace{0.55cm}%
  \begin{minipage}{0.9\textwidth}
    \vspace{1em}
    \small
    \textbf{Two-Way Tables --- Tiered Practice}\\[6pt]
    Work on the tier assigned by your teacher:
    \begin{itemize}
      \setlength{\itemsep}{0.4em}\setlength{\topsep}{0em}
      \item \textbf{Tier R:} Given a completed 2$\times$2 table --- compute conditional relative frequencies and state association.
      \item \textbf{Tier A:} Build a 2$\times$3 table from counts --- compute conditionals, sketch segmented bar chart, state association.
      \item \textbf{Tier E:} Start from joint relative frequencies --- derive conditionals, compare, state association, and investigate symmetry.
    \end{itemize}
    \vspace{0.5cm}
    {\color{redacc}\bfseries\small TIME: 10--12 minutes}
  \end{minipage}
\end{frame}
}

% ─────────────────────────────────────────────────────────────────────────────
% SLIDE 10 — EXIT TICKET
% ─────────────────────────────────────────────────────────────────────────────
{
\setbeamercolor{background canvas}{bg=navy}
\setbeamercolor{enumerate item}{fg=white}
\begin{frame}[t, plain]
  \begin{tikzpicture}[remember picture, overlay]
    \fill[navylight] (current page.north west)
      rectangle ([yshift=-1.35cm]current page.north east);
  \end{tikzpicture}
  \vspace{0.45cm}\par
  \hspace{0.55cm}%
  \begin{minipage}{0.9\textwidth}
    {\color{goldacc}\bfseries\footnotesize\MakeUppercase{Exit Ticket}}
    \vspace{0.4cm}\par
    {\color{skymid}\small\itshape
      180 students: pet ownership (Yes/No) $\times$ volunteering (Yes/No).
      Pet owners: 48 volunteer, 72 do not. Non-owners: 24 volunteer, 36 do not.}
    \vspace{0.4cm}
    \small
    \begin{enumerate}
      \setlength{\itemsep}{0.45em}
      \item {\color{white}Fill in all missing marginal totals.}
      \item {\color{white}Compute P(Volunteer $|$ Pet: Yes) and P(Volunteer $|$ Pet: No).}
      \item {\color{white}Is there evidence of an association? Justify.}
    \end{enumerate}
  \end{minipage}
\end{frame}
}

% ─────────────────────────────────────────────────────────────────────────────
% SLIDE 11 — CLOSING
% ─────────────────────────────────────────────────────────────────────────────
{
\setbeamercolor{background canvas}{bg=navy}
\setbeamercolor{itemize item}{fg=white}
\begin{frame}[t, plain]
  \begin{tikzpicture}[remember picture, overlay]
    \fill[navylight] (current page.north west)
      rectangle ([yshift=-1.35cm]current page.north east);
  \end{tikzpicture}
  \vspace{0.45cm}\par
  \hspace{0.55cm}%
  \begin{minipage}{0.9\textwidth}
    {\color{goldacc}\bfseries\footnotesize\MakeUppercase{Closing \& What's Next}}
    \vspace{0.45cm}\par
    \begin{columns}[T]
      \begin{column}{0.54\textwidth}
        \small
        {\color{goldacc}\bfseries Key Reminders}\\[5pt]
        \begin{itemize}
          \setlength{\itemsep}{0.4em}
          \item {\color{white}\textbf{Joint vs.\ conditional.}
            {\color{skymid} Joint $\div$ grand total; conditional $\div$ row/col total.
            Use conditional to detect association.}}
          \item {\color{white}\textbf{Condition on the explanatory variable.}
            {\color{skymid} Divide by the row/column of the variable that explains.}}
          \item {\color{white}\textbf{Segmented bars tell the story.}
            {\color{skymid} Same bars $=$ no association; different bars $=$ association.}}
        \end{itemize}
      \end{column}
      \begin{column}{0.42\textwidth}
        \small
        {\color{goldacc}\bfseries Where We Go Next}\\[5pt]
        \begin{itemize}
          \setlength{\itemsep}{0.35em}
          \item {\color{white}\textbf{Lesson 2.3:} Quantify association with
            difference in proportions \& relative risk}
          \item {\color{white}\textbf{Lessons 2.4--2.9:} Two quantitative variables
            $\to$ scatterplots and regression}
          \item {\color{white}\textbf{Unit 8:} Chi-square test for significance}
        \end{itemize}
        \vspace{0.5cm}
        {\color{skymid}\itshape\small Homework due next class.}
      \end{column}
    \end{columns}
  \end{minipage}
\end{frame}
}

\end{document}
